\documentclass[../../../../SoftwareRequirementsSpecification.tex]{subfiles}
\begin{document}
% Richiede le macro definite in src/macros/useCaseMacros.tex
% (incluse nel preambolo di SoftwareRequirementsSpecification.tex)

\ucstart
\begin{xltabular}{\textwidth}{|l|X|}
  \hline
  \ucid{PT-05} \\ \hline
  \ucname{Assegnazione Scheda ad un Atleta} \\ \hline
  \ucactor{PT} \\ \hline
  \ucdescription{Il PT assegna ad un atleta una scheda
  già presente nel sistema.} \\ \hline
  \ucprecond{Il PT deve essere autenticato. Deve esistere almeno una
  scheda configurata.} \\ \hline
  \ucpostcond{La scheda scelta è assegnata all'atleta e
    l'assegnazione è confermata. I carichi di tutte le sessioni della
    scheda sono specifici per quell'atleta. Se il PT ha interrotto la
    compilazione, l'assegnazione resta in bozza e non è visibile
  all'atleta (vedi Nota 5).} \\ \hline
  \ucmainflow{
  \item Il PT si trova sulla pagina di dettaglio di un atleta,
    raggiunta tramite il caso d'uso PT-08.
  \item Il PT preme il pulsante "Assegna nuova Scheda".
  \item Il sistema mostra la lista delle schede presenti nel sistema
    create dal PT.
  \item Il PT preme sulla scheda che vuole usare.
  \item Il sistema mostra il form dei dati aggiuntivi necessari: la
    data di inizio della scheda e la durata della scheda, derivata
    dal numero di blocchi mensili configurati in PT-04. Dopo che il
    PT inserisce la data di inizio, il sistema mostra anche la data
    di fine (vedi le Note).
  \item Il PT inserisce i dati e preme il pulsante "Continua".
  \item Il sistema mostra il form dei carichi degli esercizi,
    sessione per sessione, per tutte le sessioni di tutti i blocchi
    mensili della scheda (vedi le Note).
  \item Il PT inserisce i dati nel form e preme il pulsante
    "Assegna Scheda".
  \item Il sistema verifica che ogni esercizio di ogni sessione
    abbia un carico e che l'intervallo di date non si sovrapponga a
    quello di un'altra assegnazione confermata dell'atleta.
  \item Il sistema conferma l'assegnazione della Scheda all'Atleta.
  \item Il sistema invia un'email di notifica all'Atleta per
    informarlo che c'è una nuova scheda assegnata.
  } \\ \hline
  \ucaltflow{
    \textbf{4a. Creazione scheda ad hoc} \newline
    - Invece di scegliere una scheda già presente, il PT può crearne
    una da zero per l'atleta. \newline
    - Il flusso continua dal passo 1 del caso d'uso PT-03, poi dal
    passo 1 del caso d'uso PT-04. \newline
    - Dopo la creazione della scheda, il flusso riprende dal
    passo 5. \newline
    \newline
    \textbf{6a. Sovrapposizione con una scheda già assegnata} \newline
    - Il sistema rileva che l'intervallo di date della nuova
    assegnazione si sovrappone a quello di una scheda già assegnata
    all'atleta. \newline
    - Il sistema mostra un messaggio di errore con il nome e
    l'intervallo di date della scheda in conflitto. \newline
    - Il flusso riprende dal passo 5. \newline
    \newline
    \textbf{8a. Salvataggio parziale} \newline
    - Invece di premere "Assegna Scheda", il PT preme il pulsante "Salva
    e continua dopo". \newline
    - Il sistema salva l'assegnazione in bozza, con i carichi
    inseriti fino a quel momento. Non invia alcuna email all'atleta e
    non verifica la completezza dei carichi. \newline
    - Il flusso termina. Il PT riprende in un secondo momento dal
    flusso 2a. \newline
    \newline
    \textbf{2a. Ripresa di un'assegnazione in bozza} \newline
    - Sulla pagina di dettaglio dell'atleta il PT preme su
    un'assegnazione in bozza invece del pulsante "Assegna nuova
    Scheda". \newline
    - Il sistema mostra il form del passo 7 precompilato con i
    carichi già inseriti, segnalando quelli mancanti. \newline
    - Il flusso riprende dal passo 8. \newline
    \newline
    \textbf{9a. Carichi incompleti} \newline
    - Alla conferma manca il carico di almeno un esercizio. \newline
    - Il sistema segnala le sessioni incomplete e non conferma
    l'assegnazione. L'assegnazione resta in bozza. \newline
    - Il flusso riprende dal passo 8. \newline
    \newline
    \textbf{9b. Sovrapposizione rilevata alla conferma} \newline
    - Nel tempo intercorso tra la creazione della bozza e la
    conferma, all'atleta è stata assegnata un'altra scheda il cui
    intervallo si sovrappone. \newline
    - Il sistema mostra lo stesso errore del flusso 6a e non conferma
    l'assegnazione. \newline
    - Il flusso riprende dal passo 5, dove il PT può spostare la data
    di inizio. \newline
  } \\ \hline
  \ucnote{
    \textbf{Nota 1:} il PT inserisce un carico per ciascun esercizio
    di ciascuna sessione della scheda. Il sistema non calcola nessun
    carico: la variazione del peso da una sessione all'altra è
    decisa dal PT, esercizio per esercizio. Per velocizzare
    l'inserimento, il sistema precompila ogni sessione con i carichi
    della sessione precedente, che il PT conferma o modifica. Il
    campo è obbligatorio per confermare l'assegnazione, ma non per
    salvarla in bozza (vedi Nota 5).\newline
    \textbf{Nota 2:} per gli esercizi con variante drop-set, il PT
    inserisce un solo carico. Il sistema ottiene il peso delle
    ripetizioni in drop applicando il moltiplicatore configurato
    nella variante.\newline
    \textbf{Nota 3:} l'intervallo di date di un'assegnazione va dalla
    data di inizio all'ultimo giorno dell'ultimo blocco mensile. Ogni
    blocco dura un certo numero di settimane. La data di fine è quindi
    determinata dalla data di inizio e dal numero di blocchi: il
    sistema non la chiede al PT. Due assegnazioni contigue, in cui
    una termina il giorno prima che l'altra inizi, non si sovrappongono
    e sono quindi ammesse.\newline
    \textbf{Nota 4:} per adeguare i carichi di una scheda già in
    corso, il PT usa il caso d'uso PT-09, oppure assegna una nuova
    scheda.\newline
    \textbf{Nota 5:} la compilazione dei carichi può essere lunga,
    perché copre ogni esercizio di ogni sessione di ogni blocco
    mensile. Per questo il PT può interromperla e riprenderla: il
    sistema salva l'assegnazione in \textbf{bozza}. Un'assegnazione
    in bozza:
    \begin{itemize}
      \item è visibile al solo PT, che la trova nel dettaglio
        dell'atleta (caso d'uso PT-08) segnalata come tale;
      \item non è visibile all'atleta, che nel caso d'uso A-03 vede
        solo le assegnazioni confermate;
      \item non genera nessuna notifica: l'email all'atleta parte
        solo alla conferma;
      \item non è personalizzabile con il caso d'uso PT-09, che
        opera solo su assegnazioni confermate;
      \item non occupa l'intervallo di date: il controllo di
        sovrapposizione del flusso 6a considera solo le assegnazioni
        confermate. Una bozza abbandonata non impedisce quindi di
        assegnare altre schede all'atleta. Per questo motivo il
        controllo viene rifatto alla conferma (flusso 9b).
    \end{itemize}
    Lo stato di bozza riguarda solo la compilazione
    dell'assegnazione, non il periodo di allenamento. Se la scheda
    sia programmata, in corso o terminata dipende dalle sue date e
    dalla data odierna, ed è indipendente dal fatto che
    l'assegnazione sia in bozza o confermata.
  } \\ \hline
\end{xltabular}
\end{document}
